% ----------------------
\subsection{Section 9 (April 1829)}
% ----------------------
	\label{DC9}
	
	% -----
	\subsubsection{Note on the JSP and BC versions}
	% -----
	
		The version featured in the JSP comes from BC, so for this section they are the same; the BC is therefore not included in the alternate versions below.
	
	% -----
	\subsubsection{Historical Background}
	% -----
		\label{DC9-HB}
		
		% --
		\paragraph{JSP}
		% --
		
			``Earlier in April 1829, JS had dictated one revelation granting Oliver Cowdery the gift to translate ancient records and another instructing him to rely on the Holy Ghost as he translated. When Cowdery attempted to translate, however, he was unsuccessful. The revelation featured here explained Cowdery's failure and promised that he would translate at another time. Likely dictated in the second half of April, this revelation informed Cowdery that he `could have translated'' if he had proceeded correctly. He had apparently begun well but did not continue in the same manner, having `not understood'' the process: `You have supposed that I would give it unto you, when you took no thought, save it was to ask me.'' While suspending for the present Cowdery's gift to translate, the revelation underscored his ongoing current responsibility: `Behold the work which you are called to do, is to write for my servant Joseph.'''
		
	% -----
	\subsubsection{Versions}
	% -----
		\label{DC9-V}
		
		\begin{tabular}{p{1in}p{2in}p{1in}p{2in}}
			JSP		& \hyperref[EMS-1828-1829-JS-APR-1829-D]{April 1829}	& BC & Chapter 8, 20--21 \\
			DC 1835	& Section 35, 162--163									& TS & 3:18 (15 July 1842), 854 \\
			MS		& 3:8 (December 1842), 135								& DC 1844 & Section 35, 239--240 \\
		\end{tabular}
		
		\medskip
		
		See the \href{https://drive.google.com/file/d/1goAvNz8jS4R8slYfMG_db55Ty2z_vmgK/view?usp=sharing}{sections comparison} document for more information about the version history.

	% -----
	\subsubsection{Section 9:1} % *DC9-1 
	% -----
		\label{DC9-1}

			\footnote{%
				% \label{}%
				BC, which is the JSP version listed here, has the heading, ``1 \textit{A Revelation given to Oliver [Cowdery], in Harmony, Pennsylvania, April, 1829.}'' RB1 has, ``8\supers{\uline{th}.} Commandment AD 1829 A Revelation to Oliver [Cowdery] he was disrous [desirous] to know the reason why he could not Translate <\&> thus said the Lord unto him Recd. in harmony Susquehannah County Pennsylvania.'' DC 1835 has, ``\textit{Revelation given to Oliver Cowdery, April, 1829.}'' DC 1844 has, ``\textit{Revelation given to Oliver Cowdery, April, 1829.}''
				}%
		BEHOLD, I say unto you, my son, that because you did not translate according to that which you desired of me, and did commence again to write for my servant, Joseph Smith, Jun., even so I would that ye should continue until you have finished this record,
			\footnote{%
				% \label{}%
				The JSP/BC version has ``which I have intrusted unto you'' here. It's a pretty minor difference that seems easy to explain.
				}%
		which I have entrusted unto him.

		\begin{framed}
		\scriptsize
			\begin{compactdesc}
				\item[JSP] BEHOLD I say unto you, my son, that, because you did not translate according to that which you desired of me, and did commence again to write for my servant Joseph, even so I would that you should continue until you have finished this record, which I have intrusted unto you:
				\item[RB1] Behol[d] I say unto you my Son that because ye did not Translate according to that which ye desired of me \& did commence again to write for my servent Joseph even so I would that ye Should continue until ye have finished this Record which I have entrusted unto you
				\item[DC 1835] 1 Behold I say unto you, my son, that because you did not translate according to that which you desired of me, and did commence again to write for my servant Joseph Smith, jr. even so I would that you should continue until you have finished this record, which I have intrusted unto him:
				\item[DC 1844] 1 Behold I say unto you, my son, that because you did not translate according to that which you desired of me, and did commence again to write for my servant Joseph Smith, jr. even so I would that you should continue until you have finished this record, which I have intrusted unto him: 
			\end{compactdesc}
		\end{framed}

	% -----
	\subsubsection{Section 9:2} % *DC9-2 
	% -----
		\label{DC9-2}

		And then, behold, other records have I, that I will give unto you power that you may assist to translate.

		\begin{framed}
		\scriptsize
			\begin{compactdesc}
				\item[JSP] and then behold, other records have I, that I will give unto you power that you may assist to translate.
				\item[RB1] \& then Behold other Records have I that I will give unto you power that ye may assist to Translate
				\item[DC 1835] and then behold other records have I, that I will give unto you power that you may assist to translate.
				\item[DC 1844] and then behold, other records have I, that I will give unto you power that you may assist to translate.
			\end{compactdesc}
		\end{framed}

	% -----
	\subsubsection{Section 9:3} % *DC9-3 
	% -----
		\label{DC9-3}

		Be patient, my son, for it is wisdom in me, and it is not \hyperref[EXPEDIENT]{expedient} that you should translate at this present time.

		\begin{framed}
		\scriptsize
			\begin{compactdesc}
				\item[JSP] 2 Be patient my son, for it is wisdom in me, and it is not expedient that you should translate at this present time. 
				\item[RB1] be patient my Son for it is wisdom in me \& it is not expedient that ye should translate at this time
				\item[DC 1835] 2 Be patient my son, for it is wisdom in me, and it is not expedient that you should translate at this present time.
				\item[DC 1844] 2 Be patient my son, for it is wisdom in me, and it is not expedient that you should translate at this present time.
			\end{compactdesc}
		\end{framed}

	% -----
	\subsubsection{Section 9:4} % *DC9-4 
	% -----
		\label{DC9-4}

		Behold, the work which you are called to do is to write for my servant Joseph.

		\begin{framed}
		\scriptsize
			\begin{compactdesc}
				\item[JSP] Behold the work which you are called to do, is to write for my servant Joseph;
				\item[RB1] Behold this is the work which ye are called to do is to write for my Servent 
				\item[DC 1835] Behold the work which you are called to do, is to write for my servant Joseph; 
				\item[DC 1844] Behold the work which you are called to do, is to write for my servant Joseph;
			\end{compactdesc}
		\end{framed}

	% -----
	\subsubsection{Section 9:5} % *DC9-5 
	% -----
		\label{DC9-5}

		And, behold, it is because that you did not continue as you commenced, when you began to translate, that I have taken away this privilege from you.

		\begin{framed}
		\scriptsize
			\begin{compactdesc}
				\item[JSP] and behold it is because that you did not continue as you commenced, when you begun to translate, that I have taken away this privilege from you.
				\item[RB1] \& Behold it is because that ye did not continue as ye commenced when ye commenced to Translate that I have taken away this privilege from you
				\item[DC 1835] and behold it is because that you did not continue as you commenced, when you began to translate, that I have taken away this privilege from you.
				\item[DC 1844] and behold it is because that you did not continue as you commenced, when you began to translate, that I have taken away this privilege from you.
			\end{compactdesc}
		\end{framed}

	% -----
	\subsubsection{Section 9:6} % *DC9-6 
	% -----
		\label{DC9-6}

		Do not murmur, my son, for it is wisdom in me that I have dealt with you after this manner.

		\begin{framed}
		\scriptsize
			\begin{compactdesc}
				\item[JSP] Do not murmur my son, for it is wisdom in me that I have dealt with you after this manner.
				\item[RB1] do not murmer my Son for it is wisdom in me that I have dealt with you after this manner
				\item[DC 1835] Do not murmur my son, for it is wisdom in me that I have dealt with you after this manner.
				\item[DC 1844] Do not murmur my son, for it is wisdom in me that I have dealt with you after this manner.
			\end{compactdesc}
		\end{framed}

	% -----
	\subsubsection{Section 9:7} % *DC9-7 
	% -----
		\label{DC9-7}

		Behold, you have not understood; you have supposed that I would give it unto you, when you took no thought save%
			\footnote{%
				% \label{}%
				This construction here with ``save'' can be confusing. Definition \#10 of ``save'' from the 1828 says: ``To except; to reserve from a general admission or account.'' It then gives examples from \hyperref[JOSHUA-11-13]{Joshua 11:13} and \hyperref[2CORINTHIANS-11-24]{2 Corinthians 11:24}: ``Israel burned none of them, save Hazor only,'' and ``Of the Jews five times received I forty stripes, save one.'' The dictionary then has this note: ``Save is here a verb followed by an object. It is the imperative used without a specific nominative; but it is now less frequently used than except.'' So this dictionary takes the use of ``save'' like we see in verse 7 here not as a conjunction or something like that but an imperative verb which ends up having the same meaning as ``except.'' So plugging that into the verse gives you this: ``You have supposed that I would give it unto you, when you took no thought except to ask me.'' Changing the verse even more to bring out the meaning:
					\begin{quote}
						You haven't understood. You've supposed that I would give it to you even though all you did was ask.
					\end{quote}
				So the Lord is trying to communicate here that, although Oliver did ask, asking was not enough. He was required to do more than ask---and apparently, as we read in \hyperref[DC9-5]{verse 5}, he did this extra at the beginning, but then didn't keep going that way, and so his effort kind of fell apart. In the next verse the Lord explains what more was required than asking.
				}
		it was to ask me.

		\begin{framed}
		\scriptsize
			\begin{compactdesc}
				\item[JSP] 3 Behold you have not understood, you have supposed that I would give it unto you, when you took no thought, save it was to ask me;
				\item[RB1] Behold ye have not understood ye have Supposed that I would give it unto you when ye took no thought save it was to ask me 
				\item[DC 1835] 3 Behold you have not understood, you have supposed that I would give it unto you, when you took no thought, save it was to ask me;
				\item[DC 1844] 3 Behold you have not understood; you have supposed that I would give it unto you, when you took no thought, save it was to ask me; 
			\end{compactdesc}
		\end{framed}

	% -----
	\subsubsection{Section 9:8} % *DC9-8 
	% -----
		\label{DC9-8}

		But, behold, I say unto you, that you must study it out in your mind; then you must ask me if it be right, and if it is right I will cause that your \hyperref[BOSOM]{bosom} shall burn within you; therefore, you shall feel that it is right.

		\begin{framed}
		\scriptsize
			\begin{compactdesc}
				\item[JSP] but behold I say unto you, that you must study it out in your mind; then you must ask me if it be right, and if it is right, I will cause that your bosom shall burn within you: therefore, you shall feel that it is right;
				\item[RB1] but Behold I say unto you that ye must study it out in your mind then ye must ask me if it be right \& if it is right I will cause that your bosom shall burn within you therefore ye shall feel that it is right 
				\item[DC 1835] but behold I say unto you, that you must study it out in your mind; then you must ask me if it be right, and if it is right, I will cause that your bosom shall burn within you: therefore, you shall feel that it is right; 
				\item[DC 1844] but behold I say unto you, that you must study it out in your mind; then you must ask me if it be right, and if it is right, I will cause that your bosom shall burn within you; therefore, you shall feel that it is right: 
			\end{compactdesc}
		\end{framed}

	% -----
	\subsubsection{Section 9:9} % *DC9-9 
	% -----
		\label{DC9-9}

		But if it be not right you shall have no such feelings, but you shall have a stupor of thought that shall cause you to \hyperref[FORGET]{forget} the thing which is wrong; therefore, you cannot write that which is sacred save it be given you from me.

		\begin{framed}
		\scriptsize
			\begin{compactdesc}
				\item[JSP] but if it be not right, you shall have no such feelings, but you shall have a stupor of thought, that shall cause you to forget the thing which is wrong: therefore, you cannot write that which is sacred, save it be given you from me.
				\item[RB1] but if it be not right ye shall have no such feelings but ye shall have stupor of thought that shall cause you to forget the thing which is wrong there ye cannot wriete that which is sacred save it be given unto you [\textit{missing}]
				\item[DC 1835] but if it be not right, you shall have no such feelings, but you shall have a stupor of thought, that shall cause you to forget the thing which is wrong: therefore, you cannot write that which is sacred, save it be given you from me.
				\item[DC 1844] but if it be not right, you shall have no such feelings, but you shall have a stupor of thought, that shall cause you to forget the thing which is wrong: therefore, you cannot write that which is sacred, save it be given you from me.
			\end{compactdesc}
		\end{framed}

	% -----
	\subsubsection{Section 9:10} % *DC9-10 
	% -----
		\label{DC9-10}

		Now, if you had known this you could have translated; nevertheless, it is not \hyperref[EXPEDIENT]{expedient} that you should translate now.

		\begin{framed}
		\scriptsize
			\begin{compactdesc}
				\item[JSP] 4 Now if you had known this, you could have translated: nevertheless, it is not expedient that you should translate now. 
				\item[RB1] [\textit{missing}]
				\item[DC 1835] 4 Now if you had known this, you could have translated; nevertheless, it is not expedient that you should translate now. 
				\item[DC 1844] 4 Now if you had known this, you could have translated; nevertheless, it is not expedient that you should translate now.
			\end{compactdesc}
		\end{framed}

	% -----
	\subsubsection{Section 9:11} % *DC9-11 
	% -----
		\label{DC9-11}

		Behold, it was \hyperref[EXPEDIENT]{expedient} when you commenced; but you feared,%
			\footnote{%
				% \label{}%
				I think this part is crucial to seeing what Oliver's misunderstanding was, and why he had it. He did not understand a certain principle about revelation---one involving a need on his part to study out the issue in his own mind. He thought that you just had to ask without doing that ``studying out.'' It seems like he did this at the beginning (verse 5), but then stopped at some point. Why did he stop? We don't know exactly, but this verse says that he ``feared.'' I don't know exactly what that means either, but it could be that as he went along, he got scared that maybe he wasn't doing it right, or maybe he was scared that he wasn't up to the ask or that other people would mock what he had come up with. For whatever reason he ``feared''---and I take this as the explanation for why he \textit{stopped} doing things as he was doing them at the start, when he seems to have been studying it out in his mind. He started \textit{just} relying on the Lord by doing \textit{nothing but asking}. This is very important because normally we think that total reliance on the Lord is a good thing, or is a demonstration of faith, but here the Lord is trying to teach Oliver that that attitude, in this case, is the result of fear. If he'd been \textit{less} scared, and trusted the Lord \textit{more}, he would have understood that he \textit{needed} to do his part in order for the process to work correctly. But because he spooked he said, ``I'm just going to rely on the Lord, he'll give me what I need without me doing anything.'' And in this case, as it must be in many other cases as well, that was wrong. This is a very powerful lesson. Note that it also changes the way we read the previous revelations of April 1829 given to Oliver Cowdery, such as \hyperref[DC8-1]{8:1, 9--11} (what does it mean to ask in faith?), \hyperref[DC6-5]{6:5} (what does it mean to ask and knock?), \hyperref[DC6-11]{6:11} (what does it mean to inquire?), \hyperref[DC6-14]{6:14--16}, the ``fear'' in \hyperref[DC6-33]{6:33--34, 36}---I mean, Cowdery would be forgiven for thinking he just had to ask based on some of these passages. But that was not correct, or at least, it was not the lesson that Cowdery needed to learn at this particular time.
				}
		and the time is past, and it is not \hyperref[EXPEDIENT]{expedient} now;

		\begin{framed}
		\scriptsize
			\begin{compactdesc}
				\item[JSP] Behold it was expedient when you commenced, but you feared and the time is past, that it is not expedient now:
				\item[RB1] [\textit{missing}]
				\item[DC 1835] Behold it was expedient when you commenced, but you feared and the time is past, and it is not expedient now:
				\item[DC 1844] Behold it was expedient when you commenced, but you feared and the time is past, and it is not expedient now: 
			\end{compactdesc}
		\end{framed}

	% -----
	\subsubsection{Section 9:12} % *DC9-12 
	% -----
		\label{DC9-12}

		For, do you not behold that I have given unto my servant Joseph sufficient strength, whereby it is made up?  And neither of you have I condemned.

		\begin{framed}
		\scriptsize
			\begin{compactdesc}
				\item[JSP] for, do you not behold that I have given unto my servant Joseph sufficient strength, whereby it is made up? and neither of you have I condemned.
				\item[RB1] [\textit{missing}]
				\item[DC 1835] for, do you not behold that I have given unto my servant Joseph sufficient strength, whereby it is made up? and neither of you have I condemned.
				\item[DC 1844] for, do you not behold that I have given unto my servant Joseph sufficient strength, whereby it is made up? and neither of you have I condemned.
			\end{compactdesc}
		\end{framed}

	% -----
	\subsubsection{Section 9:13} % *DC9-13 
	% -----
		\label{DC9-13}

		Do this thing which I have commanded you, and you shall prosper.  Be faithful, and yield to no temptation.

		\begin{framed}
		\scriptsize
			\begin{compactdesc}
				\item[JSP] 5 Do this thing which I have commanded you, and you shall prosper. Be faithful, and yield to no temptation.
				\item[RB1] [\textit{missing}]
				\item[DC 1835] 5 Do this thing which I have commanded you, and you shall prosper. Be faithful, and yield to no temptation. 
				\item[DC 1844] 5 Do this thing which I have commanded you, and you shall prosper. Be faithful, and yield to no temptation.
			\end{compactdesc}
		\end{framed}

	% -----
	\subsubsection{Section 9:14} % *DC9-14 
	% -----
		\label{DC9-14}

		Stand fast in the work wherewith I have called you, and a hair of your head shall not be lost, and you shall be lifted up at the last day.  Amen.

		\begin{framed}
		\scriptsize
			\begin{compactdesc}
				\item[JSP] Stand fast in the work wherewith I have called you, and a hair of your head shall not be lost, and you shall be lifted up at the last day: Amen.
				\item[RB1] [\textit{missing}]
				\item[DC 1835] Stand fast in the work wherewith I have called you, and a hair of your head shall not be lost, and you shall be lifted up at the last day. Amen.
				\item[DC 1844] Stand fast in the work wherewith I have called you, and a hair of your head shall not be lost, and you shall be lifted up at the last day: Amen.
			\end{compactdesc}
		\end{framed}